\section{Введение}

Современные информационные системы всё чаще подвергаются целенаправленным атакам, использующим легитимные инструменты и обходящие традиционные средства защиты. В этих условиях возрастает роль \textbf{моделей, способных отражать логику атаки}, а не только её внешние признаки.

Традиционные системы обнаружения вторжений (IDS/IPS) в основном полагаются на сигнатурный анализ или статистические аномалии. Однако такие подходы неэффективны при столкновении с \textbf{адаптивными атаками}, использующими Living-off-the-Land (LotL) техники, где злоумышленник применяет легитимные системные утилиты (например, PowerShell, WMI, PsExec) для выполнения вредоносных действий. В таких сценариях отсутствуют чёткие сигнатуры, а поведение может не выходить за рамки статистической нормы.

В работе \cite{latypov2020multilevel} предложена \textbf{многоуровневая модель компьютерной атаки}, основанная на атрибутивных метаграфах, позволяющая формализовать взаимосвязи между целями злоумышленника, тактиками и методами реализации. Эта модель предоставляет теоретическую основу для построения семантически насыщенных систем анализа угроз. Однако в указанной работе не была детализирована \textbf{практическая методика выявления атак} на основе этой модели.

Настоящая статья устраняет этот пробел и предлагает \textbf{комплексную методику}, пригодную для внедрения в системы мониторинга безопасности, соответствующие требованиям ФСТЭК России и ГОСТ Р 56939–2024. Методика включает не только формализацию, но и алгоритмическую реализацию, оценку эффективности и рекомендации по интеграции.